\Askhsh{22174}{4}
\begin{ylh}{1}
    \item 2.1. Εξίσωση Ευθείας
    \item 3.3 Η Έλλειψη
    \item 3.4 Η Υπερβολή
\end{ylh}
\begin{wrapfigure}[15]{r}{6cm}
    \centering
    \begin{tikzpicture}[scale=0.8]
        % Define the origin
        \tkzDefPoint(0,0){O}
        % Define the foci E and E' (gamma = 3, as c=3 from problem data)
        \tkzDefPoint(3,0){E}
        \tkzDefPoint(-3,0){E'}
        
        % Define the point Gamma on the ellipse
        \tkzDefPoint(3, 16/5){Gamma} % Gamma(3, 3.2)
        
        % Draw the ellipse (major semi-axis a=5, minor semi-axis b=4)
        \draw[thick] (0,0) ellipse (5cm and 4cm);
        
        % Draw coordinate axes
        \draw[thin, gray!50] (-5.5,0) -- (5.5,0); % X-axis
        \draw[thin, gray!50] (0,-4.5) -- (0,4.5); % Y-axis
        
        % Calculate two points for drawing the tangent line 3x + 5y = 25
        % Intersection with x-axis: (25/3, 0)
        % Intersection with y-axis: (0, 5) (this is point Delta from part (b)i)
        \tkzDefPoint(25/3, 0){Px}
        \tkzDefPoint(0, 5){Py} % This is point Delta(0,5)
        
        % Draw the tangent line passing through Px and Py (and Gamma), extending slightly
        % add=-0.2 and 0.5 extends the line segment (Px, Py) by 20% before Px and 50% after Py.
        \tkzDrawLine[add=-0.2 and 0.5, thick, maincolor](Px,Py)
        
        % Draw points E, E', Gamma
        \tkzDrawPoints[fill=black, size=4pt](E,E',Gamma)
        
        % Label points
        \tkzLabelPoint[below](E){E}
        \tkzLabelPoint[left](E'){E'}
        \tkzLabelPoint[above right](Gamma){$\Gamma$}
        \node[below right, font=\footnotesize] at (E) {ήλιος};
        
    \end{tikzpicture}
\end{wrapfigure}

ΘΕΜΑ 4

Πλανήτης κινείται πάνω σε επίπεδο, ελλειπτικά γύρω από τον ήλιο του. Στο καρτεσιανό επίπεδο ο ήλιος βρίσκεται στην εστία της έλλειψης $E(\gamma,0)$, ενώ η άλλη εστία είναι στο $E'(-\gamma,0)$. Η εκκεντρότητα της τροχιάς είναι $0,6$ ενώ ο μεγάλος άξονας $10$.

\begin{alist}
    \item Να βρεθεί η εξίσωση της τροχιάς. \monades{09}
    \item Θεωρούμε ότι ο πλανήτης κινείται πάνω στην $\frac{x^2}{25} + \frac{y^2}{16} = 1$.
    \begin{rlist}
        \item Τη στιγμή που ο πλανήτης βρίσκεται στο σημείο $\Gamma\left(3,\frac{16}{5}\right)$ εκπέμπεται από αυτόν σήμα που κινείται κατά τη διεύθυνση της εφαπτομένης της τροχιάς του προς τη μεριά του άξονα $Oy$. Να εξετάσετε αν αυτό το σήμα θα περάσει από το σημείο $\Delta(0,5)$. \monades{09}
        \item Κομήτης κινείται στο ίδιο επίπεδο με τον πλανήτη και πάνω στην καμπύλη $\frac{x^2}{25} - \frac{y^2}{16} = 1$ με $x > 0$. Ποια είναι τα σημεία συνάντησης των δύο τροχιών; \monades{07}
    \end{rlist}
\end{alist}